\documentclass[12pt, onecolumn, journal]{IEEEtran}

\usepackage{amsmath, amssymb}
\usepackage{graphicx}
\usepackage{cite}

\begin{document}

\title{\LARGE Connecting the Black-Scholes-Merton Equation to the Transported Probability Density Function Method}

\author{
    \IEEEauthorblockN{Christian DiPietrantonio} \\
    \IEEEauthorblockA{ME 5311: Computational Methods to Viscous Flows\\
    \today\\
    }
}

\maketitle
\vspace{-24pt}

This paper examines the mathematical connection between the Black-Scholes-Merton equation for financial derivative pricing and the transported probability density function method used in computational fluid dynamics. In order to make this comparison, foundational knowledge of the Black-Scholes-Merton equation must be presented.

In this setting, the primary derivatives of interest are options, which are securities that give “the right to buy or sell an asset, subject to certain conditions, within a specified period of time.” The two major types of options are known as “American options” and “European options.” American options can be exercised at any time up to the date of expiration outlined on the contract, which is commonly referred to as the “expiration” or “maturity date” (last day the option can be exercised). Whereas, European options can only be exercised at the specified date of expiration. It is also important to note that the “exercise” or “strike price” is the price paid for the asset when the option is exercised. The simplest example of an option would be a contract that gives a buyer the right to purchase a single share of common stock. This example is referred to by Black and Scholes as a “call option”~\cite{BlackScholes1973}.

In deriving their valuation formula for a call option, Black and Scholes assume an idealized, frictionless market with continuous trading, no transaction costs, constant interest rates, no dividends, unrestricted short selling, and the ability to borrow or lend at the risk-free rate. The option is also assumed to be European. The methodology used by Black, Scholes, and Merton in their derivation also assumes an arbitrage-free market, which means that no trading strategy can generate riskless profits. More explicitly, an arbitrage opportunity is defined as “a (self-financing) trading strategy” that starts with zero value and terminates with a positive value. Therefore, a market is arbitrage free if there are no such arbitrage opportunities~\cite{BaxterRennie1996}.

The assumptions outlined above of Black and Scholes~\cite{BlackScholes1973} and the formalization of continuous-time arbitrage framework developed by Merton~\cite{Merton1973} lead to the derivation of the following second-order linear parabolic partial differential equation for the price $V(S,t)$ of a European call option, commonly referred to as the Black-Scholes-Merton partial differential equation:
\begin{equation}
\frac{\partial V}{\partial t} + \frac{1}{2} \sigma^2 S^2 \frac{\partial^2 V}{\partial S^2} + rS \frac{\partial V}{\partial S} - rV = 0
\end{equation}

\noindent Where $S$ is the underlying stock price, $t$ is time, $r$ is the risk-free interest rate, and $\sigma$ is the volatility of the stock, a measure of how much the stock price fluctuates or on other words “measure of our uncertainty about the returns provided by the stock”~\cite{Hull2022}. This equation was obtained by forming a hedged portfolio that eliminates risk and requiring that its returns equal re risk-free rate in an arbitrage-free market. Equation 1 is the foundation of the Black-Scholes option pricing model. 
Again, in order to advance to the mathematical connection between the Black-Scholes-Merton equation for financial derivative pricing and the transported probability density function method, a basic understanding of the transported probability density function must be obtained. For high Reynolds number flows, a statistical approach can be used to solve for the properties of the flow. In Turbulent Flows, Stephen B. Pope remarks that “in a turbulent flow, the velocity field $U(x,t)$ is random”~\cite{Pope2000}. What this means, is that if one were to repeat the same experiment multiple times under the same set of conditions, a given event may, but need not occur, making it inherently unpredictable. Although a random variable, such as velocity can not be predicted, the probability of events (e.g.\ the velocity is less than 10 m/s) can be used to create a Probability Density Function (PDF) that characterizes the random variable. Formally, the probability of any event can be determined by the Cumulative Distribution Function, and the PDF is defined to be the derivative of the CDF@\. The evolution of this Eulerian PDF of velocity in space and time can be represented by a transport equation derived by Pope from the Navier-Stokes equations. The result is the following Eulerian PDF transport equation:
\begin{equation}
\frac{\partial f}{\partial t} + V_i \frac{\partial f}{\partial x_i} = - \frac{\partial}{\partial V_i} \left[f \left \langle \frac{DU_t}{Dt} \middle| \mathbf{V} \right \rangle \right]
\end{equation}

\noindent Where $f(\mathbf{V}; \mathbf{x}, t)$ is the Eulerian PDF of velocity. Equivalently, the Lagrangian PDF transport equation, or in other words, the equation for the evolution of the conditional PDF of the particle velocity is expressed as:
\begin{equation}
\frac{\partial f^*}{\partial t} + V_i \frac{\partial f^*}{\partial x_i} = - \frac{\partial}{\partial V_i} \left[f^* a_i(\mathbf{V}, \mathbf{x}, t)\right] + \frac{1}{2} b{\left(\mathbf{x}, t\right)}^2 \frac{\partial^2 f^*}{\partial V_i \partial V_i}
\end{equation}

\noindent Where $f^*\left(\mathbf{V} \middle| \mathbf{x}, t\right)$ is the conditional PDF of the particle velocity, $a_i(\mathbf{V}, \mathbf{x}, t)$ is the drift coefficient, and $b(\mathbf{x}, t)$ is the diffusion coefficient. The drift coefficient effectively describes the direction and speed of the system if no randomness were present, and the diffusion coefficient determines the intensity of the random fluctuations induced by the vector-valued Wiener process $\mathbf{W}(t)$. To derive Equation 3, Pope expressed a generalized Langevin model for the diffusion process to describe the evolution of the velocity of a fluid particle:
\begin{equation}
d \mathbf{U^*}(t) = \mathbf{\alpha}(\mathbf{U^*}(t), \mathbf{X^*}(t), t) dt + b(\mathbf{X^*}(t), t) d \mathbf{W}
\end{equation}

he then used stochastic (more specifically Itô) calculus to obtain the Fokker-Plank equation for the joint Lagrangian PDF of the particle position and velocity, $f^{*}_L(\mathbf{V}, \mathbf{x}; t)$, and then divided the result by the marginal position PDF, $f^{*}_X (\mathbf{x}, t)$~\cite{Pope2000}. 

When introducing the Langevin equation, Pope stated that it was “originally proposed as a stochastic model for the velocity of a microscopic particle undergoing Brownian motion”~\cite{Pope2000}. This is significant to the argument of this paper, because in Stochastic Calculus for Finance II- Continuous-Time Models, by Steven E. Shreve, Shreve derives the Black-Scholes-Merton partial differential equation, starting with the price of a stock “modeled by the geometric Brownian motion”~\cite{Shreve2004} illustrated in the Stochastic Differential Equation below:
\begin{equation}
dS(t) = \alpha S(t) dt + \sigma S(t) dW(t)
\end{equation}

\noindent Therefore, one must conclude that since the Black-Scholes-Merton equation and the transported PDF equation for a particle velocity were derived from the same general form SDE, they must be effectively different versions of the same general transport equation.

The final paper will further support this connection with greater mathematical detail. This will include a deeper investigation into the stochastic calculus (specifically Itô's Lemma) used to derive either equation. The paper will then provide an evaluation of the structural equivalence of either equation through variable mapping. Once this foundation has been reinforced, the investigation will transition into evaluating the available solution methods for either transport equation. The paper will demonstrate why such methods are possible by investigating the Feynman-Kac Theorum for the Black-Scholes-Merton equation and the Kolmogorov Forward Theorum for the Fokker-Planck (transported PDF) equation and how they establish the connection between Lagrangian particle-based simulations (such as the Monte Carlo method) and Eulerian grid-based solvers (such as finite difference or finite volume schemes). Furthermore, the benefits and pitfalls of using either solution technique will be assessed for different scenarios. Finally, the paper will conclude by exploring the importance of the cross-disciplinary connection between these two equations. It will specifically investigate how the generational wealth of knowledge in computationally solving PDF equations for fluid dynamics applications can be leveraged to solve problems in the world of quantitative finance.


\newpage
\bibliographystyle{IEEEtran}
\bibliography{sources} % <-- This should match your .bib filename (without .bib)

\end{document}
